\documentclass[../main.tex]{subfiles}
\begin{document}
%==========================================TIÊU ĐỀ TẬP========================================
\begin{table}[h]
    \begin{adjustwidth}{-1cm}{}
        \begin{tabular}{>{\centering\arraybackslash}p{6cm}|>{\raggedright\arraybackslash}p{12cm}}
            \multirow{5}{6cm}
            % Cột bên trái: ÉPISODE và số tập
            {\\[-40pt]
                \centering
                {\fontsize{20pt}{20pt}\selectfont{\outlinetext[1]{eptype}{white}{\flaregothic ÉPISODE}}\\[12pt]
                    {\fontsize{72pt}{72pt}\selectfont{\outlinetext[1]{epnum}{white}{\burbank 168}}}\\[6pt]       % use for normal/special episodes
                    % {\fontsize{36pt}{36pt}\selectfont{\outlinetext[1]{epnum}{white}{\burbank SÁU}}}\\[6pt]     % use for custom episodes
                    {\fontsize{20pt}{20pt}\selectfont{\outlinetext[1]{epnum}{white}{\cafeteria (13-9)}}}\color{black}}
                \\[18pt]
                % {\fontsize{14pt}{14pt}\selectfont{\outlinetext[1]{epattr}{white}{\flaregothic [I]}}}\\
                {\flaregothic\fontsize{18pt}{18pt}\selectfont\color{horror} HORREUR}\\
                \color{black}
            }
             &
            % Dòng thứ nhất: tên tập tiếng Nhật
            {\fotskip\fontsize{18pt}{18pt}\selectfont ついてくる}                                                         \\[6pt]
            % Dòng thứ hai: tên tập tiếng Anh
             & {\cafeteria\fontsize{18pt}{18pt}\selectfont On The Tail}                                              \\[4pt]
            % Dòng thứ ba: tên tập tiếng Việt
             & {\vnmsans\fontsize{18pt}{18pt}\selectfont Bám đuổi}                                                   \\[8pt]
            % Dòng thứ tư: Nhân vật xuất hiện
             & {\caviar{\fontsize{14pt}{14pt}\selectfont Nhân vật: \emp{kuon2} \emp{game} \emp{miko2} \emp{suisei}}}
            %   \\
            % Dòng thứ năm (nếu có): Nhân vật chỉ xuất hiện trong Omake
            %  & {\abv Truyện phụ:} \emp{okayu}
            \\[8pt]
            % Dòng cuối cùng: Thời gian diễn ra của tập (thường sẽ là ngày phát hành)
             & {\sen\fontsize{16pt}{16pt}\selectfont Mốc thời gian: 31/07/2022}                                      \\[18pt]
        \end{tabular}\\[8pt]
    \end{adjustwidth}
\end{table}
%=========================================TRUYỆN CHÍNH========================================
\color{black}

\txtbx[2]{
    {\color{vol} \uline{Tóm tắt:}} Lỡ chân ngồi bẹp chiếc kính râm viền vàng mang con số ``2002'' quý giá của Sakura Miko, Hoshimachi Suisei vô tình châm ngòi cho một cơn thịnh nộ long trời lở đất khiến nàng Vu nữ hóa thành ``quái thú'' điên cuồng truy sát. Cuộc đào tẩu nghẹt thở cuốn theo cả anh Tổng Kuon băng qua ma trận ngõ ngách, cánh rừng mù sương, sông nước và một hội chợ mùa hè ma quái không bóng người. Bị tách khỏi Kuon và đối mặt với đoàn tàu ngầm bí ẩn lao tới xé toạc màn đêm, Suisei trải qua cơn ác mộng kinh hoàng chỉ để rồi choàng tỉnh trên chiếc ghế sofa quen thuộc tại văn phòng, nơi một mảnh ghép thực tại kỳ lạ đang chờ đón.
}

\pov{Hoshimachi Suisei}

Một bầu không khí nặng trĩu và ngột ngạt đến nghẹt thở đang bao trùm khắp bốn góc văn phòng.

Tôi vừa mới gây ra một chuyện - có thể nói là tày đình - đối với cô bạn thân chí cốt của mình, Sakura Miko. Nói cho chính xác hơn thì tai họa không giáng trực tiếp lên người cô ấy, mà là lên một món đồ sở hữu của cô.

Đứng tựa lưng bên chiếc bàn làm việc, ngay cả anh Tổng của chúng tôi cũng chỉ biết lắc đầu ngao ngán buông một câu: ``Thôi rồi\dots\ thế này thì đúng là cơm toi thật rồi.''

Thậm chí từ mặt chiếc đồng hồ trên cổ tay anh, vị trợ lý ảo Game cũng cất giọng tặc lưỡi châm chọc: ``Vụ này thì e rằng có chất tiền tấn ra đền cũng chẳng thể nào mua chuộc nổi đâu\dots''

Trong khi đó, ở phía đối diện, khuôn mặt của Miko-chi đã hoàn toàn xây xẩm lại, tối sầm như bầu trời trước cơn giông bão. Cô ấy tức giận đến mức không thốt nên lời, hai bờ vai run lên bần bật, trán nổi đầy gân xanh, trông chẳng khác nào một con quái thú thời tiền sử đang chuẩn bị thức tỉnh để xé toạc con mồi ra làm trăm mảnh.

Tôi sợ đến toát mồ hôi hột, vội vã chắp hai tay trước ngực cúi đầu rối rít: ``Tôi xin lỗi mà! Tôi thật sự xin lỗi bà nhiều lắm, Miko-chi!''

Anh Kuon đứng bên cạnh khẽ thở dài, bước lên một bước làm cầu nối hòa hoãn: ``Thôi nào Miko-chan, bình tĩnh lại chút đi em. Dù sao thì em ấy cũng đâu có cố ý làm thế đâu\dots''

Phải nhờ đến uy tín và giọng điệu mềm mỏng của anh ấy đứng ra xoa dịu, ngọn lửa thịnh nộ trong mắt nàng Vu nữ mới tạm thời lắng xuống đôi chút để chịu mở miệng nói chuyện.

Thực ra thì\dots\ đầu đuôi câu chuyện cũng chẳng có gì to tát: Vừa nãy lúc bước vào phòng, tôi lơ đễnh ngồi phịch xuống ghế sofa mà không để ý, vô tình đè bẹp dí chiếc kính râm của cô ấy.

\img{13-9a}

Đúng vậy, sự tình quả thực chỉ có vỏn vẹn ngần ấy thôi.

``\textit{Nhưng mà Sui-chan ơi, chỉ vì một chiếc kính râm cỏn con mà Miko-chi lại nổi trận lôi đình đến mức đó sao?}'' tôi biết chắc chắn các bạn sẽ thắc mắc như vậy. Khổ nỗi, vấn đề nằm ở chỗ cái kính râm đó hoàn toàn không phải là một món đồ phụ kiện bình thường!

Anh Tổng khoanh tay, nhướng mày hỏi với vẻ vừa không đồng tình nhưng cũng xen lẫn chút tò mò: ``Em có vẻ trân quý cái kính râm ấy ghê nhỉ, Miko-chan?''

Cô bạn tôi run rẩy nhặt chiếc kính râm đã bị vỡ nát tròng kính lên, giọng nghẹn ngào trong sự thất vọng tột cùng: ``Em cũng không rõ nữa\dots\ nhưng mà cái loại kính này mang kiểu dáng thiết kế độc nhất vô nhị chỉ được sản xuất và bày bán vào đúng những năm 2002 thôi đấy ne!''

\img{13-9b}

Đó là một chiếc kính râm với phần gọng kim loại mạ vàng lấp lánh, nhưng điểm kỳ quái là hai bên mắt kính lại được uốn thành hai chữ số mười màu kim loại đặt đối xứng, ghép lại thành con số 2002 vô cùng dị hợm.

Tôi nuốt nước bọt, lo âu đáp khẽ: ``Ừ\dots\ có lẽ là thế thật\dots''

Kuon thở dài gãi trán: ``Anh vốn định bảo `chỉ là cái kính râm thôi mà sao em cứ phải giãy nảy hết cả lên', nhưng mà nghe em giải thích nguồn gốc cổ lỗ sĩ như thế thì\dots\ anh cũng hiểu được là nó quý giá tới cỡ nào rồi.''

Miko ôm chiếc kính vỡ vào lòng, nước mắt lưng tròng: ``Vĩnh biệt nhé, tuổi thơ tươi đẹp của tôi\dots''

Giọng Game bỗng đều đều cất lên từ loa đồng hồ: ``Tôi ghét phải dội một gáo nước lạnh vào lúc này, nhưng theo dữ liệu nhân khẩu học thì vào năm 2002, cô thậm chí còn chưa sinh ra trên cõi đời này nữa kìa.''

Tôi gật gù đồng tình: ``Công nhận thật.''

Ông ta lại tiếp tục bồi thêm một đòn chí mạng: ``Kể cả nếu cô có sinh ra vào đúng thời điểm đó đi chăng nữa, thì lúc ấy cô cũng còn quá bé để có thể nhận thức được bất kỳ ký ức tuổi thơ nào gắn liền với--''

``\textbf{Giờ còn quan trọng quái gì mấy cái chuyện đó nữa ne!}'' chưa kịp để vị trợ lý ảo dứt câu, Miko-chi đã quay ngoắt người lại, gào toáng lên át cả tiếng loa.

Anh Tổng ngán ngẩm khoát tay: ``Người ta nói hoàn toàn đúng thực tế mà em còn cãi cái nỗi gì. Với lại ấy\dots''

Tôi cũng tức mình chêm vào: ``Nếu nó thực sự quý giá và quan trọng với bà như vậy, thì việc quái gì bà lại vứt lăn lóc nó ngay trên ghế sofa phòng khách chứ hả?!''

Kuon tiếp lời: ``Đúng đấy. Nếu em không bất cẩn vứt bừa bãi như thế thì giờ đã chẳng phải ngồi đây khóc lóc ỉ ôi.''

Nàng Vu nữ vẫn ngoan cố cự cãi, hai mắt rớm lệ, ngón tay run rẩy chỉ thẳng về phía chiếc đệm sofa: ``Em không hề vứt nó! Em rõ ràng là để nó ở đó mà ne!''

Tôi và anh Kuon nhìn nhau, cùng lúc thốt lên với sự bất lực chạm đáy:  ``Ý bà là sao cơ chứ?!'' / ``Để với chả vứt, có khác đếch gì nhau đâu em?''

*RẮC!*

Một tiếng gãy giòn tan chói tai vang lên. Bất thình lình, Miko-chi dùng cả hai bàn tay bẻ gập đôi chiếc kính râm quý giá ấy làm hai nửa. Cô ngẩng phắt đầu lên nhìn thẳng vào chúng tôi, toàn thân tỏa ra một luồng hắc khí tím sẫm ngùn ngụt như một con ác quỷ vừa bò lên từ địa ngục. Đôi mắt trắng bệch gườm gườm căm phẫn, cô nghiến răng gằn từng chữ: ``Bà\dots\ có vẻ thích chọc ghẹo tôi lắm đúng không? Vậy giờ bà muốn ăn roi\dots\ hay là muốn ăn đấm đây ne?!''

\img{13-9c}

Hai chúng tôi cùng đơ người ra trước câu hỏi kỳ quặc đến mức khó đỡ ấy: ``Hả\dots? Chẳng phải cả hai cái đó\dots\ đều cùng là bị đánh đòn hay sao?!''

Thôi xong rồi\dots\ Không ổn chút nào rồi! Miko-chi đã thực sự phát điên vì mất đi chiếc kính đó!

Cô ấy gầm lên một tiếng man dại rồi nhảy bổ về phía chúng tôi như một con thú dữ săn mồi. Trong cơn hoảng loạn tột cùng, tôi lập tức quay phắt người lao ra phía cửa thoát hiểm, bàn tay theo phản xạ vô tình nắm chặt lấy cổ tay của anh Tổng kéo giật đi.

Tôi hét toáng lên: ``\dots Chết tiệt! Chạy mau thôi!''

Kuon bị giật ngửa ra sau, loạng choạng chạy theo: ``Này! Sao tự dưng em lại lôi cả anh đi theo thế!?''

Tiếng gào thét của Miko vang dội như sấm truyền ngay sau gáy: ``\textbf{Đừng có hòng chạy trốn khỏi tội lỗi của mình neeeeee!}''

Game buông một câu ráo hoảnh: ``Xác nhận: Đối tượng đã hoàn toàn rơi vào trạng thái mất trí cục bộ\dots''

Thế là, cuộc đào tẩu nghẹt thở chính thức bắt đầu.

\ct

Vừa cắm đầu cắm cổ chạy thục mạng ra ngoài đường lớn, tay tôi vừa siết chặt lấy cổ tay Kuon-senpai anh Kuon lôi đi xềnh xệch, miệng vừa hét lớn đáp trả lại phía sau: ``\textbf{Nếu nó quan trọng đến thế thì bà phải cất vào hộp đàng hoàng chứ! Ai bảo bà để lung tung hả!}''

\img{13-9d}

``\textbf{Này nhá!! Đứng lại đó cho tôiiii!}''

\dots Khỉ thật, cô ấy bám đuổi dai dẳng ngay phía sau lưng chúng tôi kìa!

Kuon vừa chạy vừa thở dốc, quay đầu lại phân trần: ``Em ấy nói đúng rồi còn kêu ca cái gì nữa-- C-Cái quái gì thế kia?!''

``\textbf{Đứng lại đó, neeeeeeeeeeeee!!!}''

Tôi thở hồng hộc: ``Cô ấy đúng là cứng đầu một cách điên rồ thật đấy\dots''

Game chêm vào: ``Hôm nay não bộ của cô ta bị chập mạch ở vi mạch nào vậy trời?''

Chúng tôi tiếp tục guồng chân lao về phía trước, thi thoảng ngoái đầu lại nhìn. Cứ mỗi lần tôi quay lại, hình ảnh Miko vung vẩy chiếc gậy trừ tà đuổi bén gót lại hiện ra mồn một, trông chẳng khác nào một hung thần đòi nợ.

Bình thường thì mấy trò đuổi bắt này diễn ra như cơm bữa ở công ty, thế nhưng tại sao hôm nay cô ấy lại hóa rồ đến cái mức độ đáng sợ này cơ chứ?!

Nhác thấy một ngã rẽ khuất ngay phía trước ngã tư, tôi liền hô lớn: ``Kuon-senpai, bên này!''

Tôi kéo mạnh người anh ấy để bắt kịp nhịp chạy, rồi cả hai cùng quặt gấp sang con đường nhánh bên trái.

\img{13-9e}

Một đoạn đường dài trôi qua\dots

Tôi vừa thở dốc vừa liếc mắt nhìn lại: ``Cậu ta\dots\ vẫn còn đuổi theo chúng ta sao?''

Kuon ôm ngực, thở hộc hộc từng ngụm khí: ``Vẫn còn! Cách chúng ta có một đoạn ngắn thôi đấy!''

``Điên thật chứ!''

Đến ngã tư tiếp theo, chúng tôi lại rẽ ngoặt sang bên phải, băng qua những dãy nhà san sát và ma trận hẻm hóc chằng chịt của thành phố. Dưới sự dẫn dắt của tôi, bóng dáng của cô bạn Vu nữ điên rồ kia dường như đã dần dần bị bỏ lại phía sau những bức tường gạch\dots

\img{13-9f}

Game cất giọng thông báo: ``\dots Dữ liệu nhiệt độ môi trường cho thấy: Sakura-san dường như đã mất dấu chúng ta rồi.''

Kuon chống hai tay lên đầu gối, thở không ra hơi: ``Thế\dots\ thế á\dots?''

Tôi cũng kiệt sức tựa lưng vào tường: ``Tìm\dots\ tìm chỗ nào đó nghỉ chút đi\dots\ em hết hơi rồi\dots''

Vị trợ lý ảo cảnh báo: ``Tôi nghĩ nán lại đây không phải là một ý kiến hay ho đâu\dots''

Trước mặt chúng tôi là một con ngõ cụt hẹp và vắng vẻ. Tôi lảo đảo bước vào trong, tựa cả tấm lưng ướt đẫm mồ hôi vào bức tường gạch thô nhám, lồng ngực phập phồng dữ dội. Ơn trời là thể lực và sức bền của tôi vẫn còn đủ để chống đỡ\dots

Thế nhưng, nhìn sang bên cạnh, người đồng hành của tôi thì đã hoàn toàn kiệt quệ. Anh Kuon ngã vật xuống nền đất, cả người nằm xoài ra, hai mắt nhắm nghiền, dùng cả cánh tay áo che ngang mặt để trốn tránh ánh sáng. Tôi quên khuấy mất rằng một người làm công việc văn phòng như anh ấy vốn dĩ có thể lực và sức bền khá kém. Dù vậy, trong lúc kiệt sức nhất, anh ấy vẫn có ý tứ dùng tay che mặt để giữ kẽ.

\img{13-9f1}

Kuon thều thào qua kẽ môi: ``Chưa bao giờ\dots\ trong đời\dots\ anh phải chạy\dots\ thục mạng nhiều đến cái mức này\dots''

Game tranh thủ châm chọc: ``Giá như thường ngày cậu chịu khó tập luyện thể thao một chút thì giờ đâu đến nông nỗi này, Kuon-user.''

Anh Tổng gắt khẽ trong hơi thở đứt quãng: ``Thời gian quái đâu ra mà tập với tành\dots''

Tôi bật cười trừ, đưa tay ra trước mặt anh: ``Game-san nói đúng đấy. Nào, đưa tay đây cho em, để em kéo anh dậy.''

Dùng hết chút sức lực còn sót lại, tôi nắm lấy bàn tay to lớn của anh Tổng rồi kéo mạnh một cái để anh gượng đứng dậy.

``Cảm ơn em nhé\dots'' - anh ấy đáp lại bằng giọng mệt mỏi, rồi cũng đứng sát cạnh tôi dựa lưng vào bờ tường gạch.

Cả hai chúng tôi đứng thở dốc suốt một hồi lâu để hồi lại sức. Mồ hôi thi nhau túa ra như tắm, chẳng rõ là vì vận động quá sức hay vì nỗi sợ hãi tột cùng trước sự truy đuổi điên cuồng của Miko-chi. Mà công nhận\dots\ sức bền của con nhỏ ngốc nghếch đó trâu bò thật đấy chứ.

Bỗng nhiên\dots

*LENG KENG\dots\ LENG KENG\dots*

Tim tôi như ngừng đập một nhịp. Tiếng chuông\dots\ đó chính xác là tiếng chuông bạc đeo trên cổ áo của cô ấy.

Tôi hốt hoảng: ``Cô ấy\dots\ cô ấy vẫn đang bám theo chúng ta kìa!''

Game chép miệng cạn lời: ``Thật sự luôn đấy, cô ta bám dai như đỉa đói vậy.''

Kuon ôm trán rên rỉ: ``Lại\dots\ lại phải chạy tiếp nữa à?!''

``Có vẻ là vậy rồi! Đi thôi, Kuon-senpai!''

Không chút chần chừ, tôi lại nắm chặt lấy tay anh Tổng kéo lao ra khỏi ngõ ngách. Tôi quay sang hỏi dồn: ``Anh vẫn còn chạy nổi chứ?!''

Dù mặt mũi đã cắt không còn một giọt máu, lồng ngực như muốn nổ tung, nhưng anh Tổng vẫn kiên cường giơ một ngón tay cái lên làm dấu hiệu đồng ý. Tôi thầm thấy có lỗi vô cùng: tự dưng vì chuyện của tôi mà kéo anh vào cái cuộc rượt đuổi sinh tử quái đản này\dots\ Thế nhưng, nếu như tôi không kéo anh ấy đi cùng, ai mà biết được con quái thú Miko-chi trong cơn điên loạn sẽ làm ra chuyện khủng khiếp gì với anh ấy cơ chứ?!

Thế nên, chúng tôi không có lựa chọn nào khác. Chúng tôi buộc phải tiếp tục chạy.

\ct

Băng qua thêm không biết bao nhiêu ngã rẽ và những con đường ngoằn ngoèo, dựa vào những chỉ dẫn mò mẫm của Game, cuối cùng chúng tôi cũng chạy thoát ra khỏi khu vực nội đô sầm uất.

\img{13-9g}

Thế nhưng, Miko-chi dường như vẫn không hề có ý định buông tha.

Trước mặt chúng tôi là một cánh rừng rậm rạp hoang vu. Chúng tôi lao thẳng vào màn cây cối rậm rạp, những tán lá thông sắc nhọn quẹt vào da thịt, sương mù dày đặc phả vào mặt buốt lạnh đến thấu xương\dots

\img{13-9h}

Tôi vừa chạy vừa ngoái đầu: ``Miko-chi vẫn còn đuổi theo sau chứ, senpai?''

Kuon thở hồng hộc: ``Anh\dots\ anh vẫn còn nghe văng vẳng tiếng hét thất thanh của em ấy phía sau lưng kia kìa\dots''

``Ta phải chạy nhanh hơn nữa thôi!''

Chúng tôi cắm đầu cắm cổ lao về phía trước, cho đến khi nhìn thấy một vệt sáng mờ ảo le lói ở bìa rừng\dots

Thoát ra khỏi rừng cây, cảnh tượng trước mắt lại biến đổi một cách kỳ dị: Chúng tôi đang chạy băng băng ngay trên mặt một con sông rộng lớn. Nước sông lạnh buốt bắn tung tóe dưới gót giày, trong khi Kuon-senpai vẫn đang cố hết sức bám sát từng bước chân của tôi\dots

\img{13-9i}

``\textbf{Đứng lại đó cho tôi, neeeeee!!!!}''

Tiếng thét xé toạc màn đêm của Miko-chi lại vang dội từ phía sau. Khốn thật chứ, rốt cuộc thì cô ta làm bằng cái thứ sắt đá gì mà dai dẳng đến mức này cơ chứ?!

Càng dốc sức chạy, đôi tai tôi lại càng nghe rõ mồn một những tiếng kim loại va đập vào nhau lanh lảnh. Thế nhưng\dots\ đó hoàn toàn không phải là tiếng chuông đeo cổ của ``quái thú'' Miko, mà là\dots

\dots những tiếng kẻng rộn rã phát ra từ một hội chợ mùa hè.

Tôi phanh gấp bước chân lại.

Quang cảnh xung quanh tôi đã hoàn toàn biến đổi tự lúc nào: Mặt sông nước biến mất, thay vào đó là những dãy lều bạt đỏ rực của các quầy hàng lưu niệm, cùng những cánh cổng torii đỏ thẫm dựng san sát nhau trải dài hun hút.

\img{13-9j}

Một luồng khí lạnh buốt bất chợt chạy dọc sống lưng tôi.

Tôi hoảng hốt ngoái đầu nhìn lại con đường sông vừa chạy qua: Nơi ấy giờ đây cũng đã biến thành một quầy hàng bán mặt nạ ma quái từ bao giờ. Thế nhưng, điều khiến toàn bộ tế bào trong người tôi đông cứng lại vì kinh hãi tột cùng, đó là\dots

``Kuon-senpai? Anh đâu rồi?! \textbf{Kuon-senpai ơi?!}''

Không có ai trả lời tôi cả.

Anh ấy không có ở đây! Chết tiệt thật\dots\ Chẳng lẽ vì quá kiệt sức mà anh ấy đã không thể đuổi kịp bước chân tôi và bị lạc mất rồi sao?!

Bây giờ\dots\ chỉ còn lại duy nhất một mình tôi trơ trọi giữa cái chốn quỷ quái, ma mị này.

Những tiếng kẻng kim loại cứ thế tiếp tục văng vẳng, xoay mòng mòng quanh màng nhĩ tôi, mỗi lúc một dồn dập, mỗi lúc một lớn dần như muốn xé toạc đầu óc tôi ra làm trăm mảnh. Tôi run lên bần bật trong nỗi sợ hãi tột cùng. Xung quanh tôi, ngoài những quầy hàng rực rỡ không một bóng người bán kẻ mua, không gian tĩnh lặng đến rợn người. Đó là một trải nghiệm cô độc và đáng sợ đến nghẹt thở.

Tôi vội vã đưa hai tay lên bịt chặt lấy hai tai, đầu óc hoảng loạn cực độ, có cảm giác như một thực thể tà ác vô hình nào đó đang chực chờ nuốt chửng lấy tôi.

\img{13-9k}

Tôi buột miệng chửi thề: ``Khỉ thật chứ!''

Không thể chần chừ thêm nữa, tôi nhắm nghiền mắt, cắm đầu cắm cổ lao thục mạng về phía những cánh cổng torii sâu hun hút kia. Nếu còn đứng lại nơi này thêm một giây phút nào nữa, khéo tôi cũng sẽ hóa điên mất.

Hai tay vẫn ghì chặt lấy tai để chặn đứng những âm thanh ma quái, tôi cứ thế chạy như một kẻ mộng du, không màng đến bất kỳ chướng ngại vật nào dưới chân.

Và rồi\dots

*BỊCH!*

Chân tôi vấp phải một vật cứng chắn ngang đường, khiến toàn bộ thân người tôi mất thăng bằng, ngã nhào xuống nền đất lạnh buốt.

Đùng một cái, mọi âm thanh của hội chợ mùa hè bỗng nhiên im bặt hoàn toàn. Không gian rơi vào một khoảng lặng chết chóc.

Tôi lồm cồm bò dậy, hai mắt trợn trừng khi nhận ra trước mặt mình lúc này chỉ là một đôi đường ray xe lửa bằng thép xám xịt chạy dài vô tận.

Xung quanh tôi là một màn đêm đen kịt, sâu thăm thẳm. Không một bóng cây, không một mái nhà, không có bất kỳ dấu hiệu nào của sự sống. Chỉ có những đốm đèn tín hiệu mờ ảo hắt ánh sáng le lói xuống hai thanh ray lạnh lẽo. Rất có thể\dots\ đây là một đoạn đường ray tàu điện ngầm dưới lòng đất.

\img{13-9l}

\textit{VÙ\dots}

Một luồng gió lạnh thốc mạnh vào gáy tôi từ phía sau. Khi tôi quay đầu lại, cảnh tượng đập vào mắt khiến tôi rụng rời toàn bộ tay chân\dots

Một đốm sáng trắng chói lòa bỗng lóe lên từ nơi sâu thẳm của đường hầm, kèm theo đó là tiếng kim loại nghiến ken két đinh tai nhức óc của những bánh xe sắt nặng hàng trăm tấn.

*TU TU\dots! TUUUUUU\dots!*

\img{13-9m}

Tiếng còi tàu rít lên từng hồi dài ghê rợn, xé toạc bầu không khí tĩnh mịch. Ánh đèn pha của con tàu cứ mỗi giây lại một sáng rực lên, rọi thẳng vào người tôi. Con tàu ma quái kia đang lao vun vút tới với vận tốc xé gió và hoàn toàn không có bất kỳ dấu hiệu nào cho thấy nó sẽ dừng lại.

``Cái gì thế này\dots? Tại sao lại ở đây chứ?! Mình đang ở cái chốn quái quỷ nào thế này?!'' hàng ngàn câu hỏi gào thét trong tâm trí tôi, nhưng chẳng có lấy một ai đáp lời.

Bản năng sinh tồn mách bảo tôi phải đứng dậy. Tôi thở dốc, lảo đảo gượng hai chân đứng lên và cố hết sức chạy thục mạng dọc theo con đường ray vô định ấy.

Thế nhưng, cơ thể tôi đã chạm tới giới hạn kiệt cùng. Những bước chạy của tôi trở nên loạng choạng, xiêu vẹo, đôi chân nặng trĩu như đeo chì. Trong đầu tôi chỉ găm đúng một ý nghĩ duy nhất: ``\textit{Phải chạy\dots\ phải chạy trốn khỏi cô ta\dots}'' Nhưng cả sức lực lẫn tinh thần của tôi lúc này đều đã cạn kiệt hoàn toàn.

Chẳng mấy chốc, con quái vật sắt thép khổng lồ kia đã áp sát ngay sau lưng tôi.

Không\dots\ Không thể nào\dots\ Tại sao mọi chuyện lại thành ra thế này?!

Khoảnh khắc tôi ngoảnh đầu lại nhìn, luồng ánh sáng trắng lóa mắt của ngọn đèn pha đã nuốt trọn lấy toàn bộ tầm nhìn của tôi. Con tàu lao thẳng vào tôi mà không hề giảm tốc độ.

Và rồi\dots\ điều gì phải đến cuối cùng cũng đã đến.

Tất cả những gì còn đọng lại sau đó chỉ là một màn sương trắng xóa bao trùm, hòa lẫn cùng tiếng hét thất thanh kinh hoàng của chính bản thân tôi\dots

\img{13-9n}

\dots

\dots

Nhưng\dots

Khoan đã\dots

Tôi vẫn đang nhắm nghiền hai mắt, cả người co rúm lại chờ đợi sự va chạm khủng khiếp sẽ nghiền nát thân xác mình.

Thế nhưng\dots\ tôi chẳng cảm thấy bất kỳ cơn đau đớn nào cả. Không có tiếng kim loại va đập, không có chấn động rung chuyển.

Tôi run rẩy hé mở một bên mắt.

Trước mặt tôi vẫn chỉ là một đoạn đường ray hoang vắng trong bóng tối. Con tàu khổng lồ hung hãn ban nãy\dots\ đã biến mất không một dấu vết.

Chuyện quái quỷ gì vừa mới xảy ra với tôi vậy?!

Tất cả những thứ kỳ quái vừa rồi\dots\ từ khu rừng mù sương, dòng sông nước chảy, rồi đến hội chợ mùa hè rợn người\dots\ và cả đoàn tàu điện ngầm này nữa\dots

Rốt cuộc tất cả những thứ này là cái gì đây?!

Có ai không?! Làm ơn hãy giải thích cho tôi với! Có ai ở đây không\dots\ mau đưa tôi rời khỏi cái nơi chết tiệt này đi!

Đó là những lời mà tôi muốn gào thét lên thật to, thế nhưng cổ họng tôi nghẹn đắng lại, lồng ngực không còn chút hơi tàn để bật ra thành tiếng. Sự bất lực và nỗi sợ hãi tột cùng bóp nghẹt lấy tâm trí tôi khi nhận ra mình đang hoàn toàn đơn độc.

Trong lúc tôi vẫn còn đang ngồi bệt dưới đất, toàn thân run rẩy bần bật nhìn chằm chằm vào khoảng không gian nơi đoàn tàu ma vừa biến mất\dots

\dots từ trong bóng tối phía sau lưng tôi, bất thình lình thò ra hai cánh tay.

\img{13-9o}

Không để tôi kịp định thần hay quay đầu lại, hai bàn tay ấy đã chộp lấy tôi, đi kèm với một giọng cười the the rùng rợn vang lên ngay sát vành tai:

``He he\dots\ Cuối cùng thì tôi cũng đã bắt được bà rồi ne\~{}''

Khoan đã\dots\ giọng nói này\dots\ không lẽ nào lại là\dots?!

*\textit{ROẠP!}*

Những gì diễn ra sau khoảnh khắc kinh hoàng ấy\dots\ tôi hoàn toàn không còn nhớ nổi nữa.

Tôi chỉ lờ mờ nhận thức được một điều duy nhất\dots

Con ``quái thú'' điên loạn ấy đã tóm chặt lấy tôi. Cuộc rượt đuổi đẫm nước mắt này\dots\ rốt cuộc cũng đã đi tới hồi kết\dots

%==========================================TRUYỆN PHỤ=========================================
\omk

\begin{center}
    \uline{Thực tại.\\
        Văn phòng \hll.}
\end{center}

``\textbf{Á Á Á Á Á Á Á Á!}''

Tôi hét toáng lên một tiếng thất thanh rồi bật phắt người ngồi dậy trên chiếc đệm sofa êm ái, lồng ngực phập phồng dữ dội, mồ hôi đầm đìa túa ra ướt đẫm cả vạt áo.

\img{13-9p}

Hoang mang đảo mắt nhìn quanh, tôi bàng hoàng nhận ra mình đang ngồi ngay ngắn giữa phòng tiếp khách quen thuộc của văn phòng \hll. Ánh nắng chiều dịu nhẹ xuyên qua khung cửa kính, tiếng điều hòa nhiệt độ rì rầm phả hơi mát.

Hóa ra\dots\ tất cả những điều khủng khiếp vừa rồi chỉ là một giấc mơ. Nói cho chính xác hơn thì đó là một cơn ác mộng kinh hoàng đến điếng người.

Liếc mắt sang chiếc bàn kính bên phải, tôi thấy Miko và anh Kuon đang ngồi song song trên hai chiếc ghế đơn, tay cắm cúi lướt màn hình điện thoại. Nghe thấy tiếng hét xé họng của tôi, cả hai người họ đều giật bắn mình ngẩng đầu lên, nhìn tôi bằng ánh mắt tròn xoe ngập tràn sự kinh ngạc. Trên gương mặt cô bạn Vu nữ thậm chí còn thoáng hiện lên vẻ âu lo thực sự.

Cô bạn Vu nữ chớp mắt lo lắng: ``B-Bà có sao không đấy, ne?''

\img{13-9q}

Kuon buông điện thoại xuống bàn, nhíu mày hỏi: ``Có chuyện gì thế Sui-chan? Tự dưng em lại hét toáng lên như gặp ma thế?''

Đáp lại sự khó hiểu của họ chỉ là một tiếng thở dốc đầy ngơ ngác của tôi: ``Hả\dots?''

Đầu óc tôi lúc này vẫn đang quay cuồng, hoàn toàn chưa thể phân biệt nổi đâu là thực tại, đâu là ảo ảnh mộng mị.

Từ mặt đồng hồ trên tay anh Tổng, Game cất giọng châm chọc: ``Hả với chả hở cái gì. Nãy giờ tôi thấy cô nằm ngủ mà cứ rên rỉ ú ớ suốt thôi đấy.''

Anh Kuon gật đầu phụ họa: ``Thậm chí hai tay hai chân em còn huơ loạn xạ cả lên kia kìa. Trông vất vả ghê gớm.''

Miko nghiêng đầu đoán mò: ``Chắc là bà vừa mới gặp phải một cơn ác mộng tồi tệ lắm đúng không ne\dots''

Quả đúng như tôi vừa tự trấn an bản thân, tất cả những gì tôi vừa trải qua - từ việc giẫm nát kính, Miko hóa quỷ truy sát, cho đến cánh rừng, dòng sông, hội chợ ma và đoàn tàu điện ngầm rùng rợn - hoàn toàn chỉ là một giấc chiêm bao hư ảo. Một giấc mơ không có thật.

Tôi đưa cả hai tay lên ôm lấy cái đầu vẫn còn đang ong ong nhức buốt. Cảm giác hụt hẫng và chông chênh lúc này thật kỳ lạ, chẳng khác nào vừa bước xuống sau một chuyến tàu lượn siêu tốc tử thần.

Tôi thở phào một hơi nhẹ nhõm, chống tay toan đứng dậy đi về phía cửa sổ để hóng chút gió trời cho tỉnh táo.

Thế nhưng, ngay khoảnh khắc gót chân tôi vừa hạ xuống nền nhà\dots

\textit{RẮC!}

Một tiếng nứt vỡ khô khốc, đanh gọn vang lên ngay dưới chân tôi.

Cả người tôi đông cứng lại. Tôi run rẩy cúi đầu nhìn xuống khe đệm sofa nơi mình vừa nằm ban nãy.

Ở đó\dots\ nằm chỏng trơ ngay dưới mép đệm\dots\ là một chiếc kính râm gọng kim loại mạ vàng, với hai bên mắt kính được uốn thành cặp số mười đối xứng tạo thành con số 2002. Và trên mặt tròng kính sẫm màu ấy\dots\ đang hằn rõ những vết nứt chằng chịt.

Máu trong người tôi như đông cứng lại. Mồ hôi lạnh lập tức túa ra ướt đẫm sống lưng.

Không\dots\ Không lẽ nào\dots\ Tôi vẫn chưa thực sự thoát khỏi cơn ác mộng đó sao?! Hay là\dots\ tất cả những gì tôi nhìn thấy trong giấc mơ ban nãy\dots\ thực chất chính là một sự báo trước về tương lai?!

``\dots Sui-chan? Em sao thế, Sui-chan?''

``Hoshimachi-san?''

``Này Sui-chan ơi!''

Tiếng gọi dồn dập của cả ba người họ kéo phăng tôi trở về với thực tại. Tôi giật mình thon thót, mặt mày tái mét không còn một giọt máu, lắp bắp không nên lời: ``D-Dạ\dots\ thì là\dots\ em\dots''

Đúng lúc tôi đang bối rối tột cùng không biết phải giải thích ra sao, thì Miko-chi bỗng thốt lên: ``Ủa, cái gì ở dưới chân bà thế kia ne?''

Cô bạn Vu nữ nhanh nhẹn tiến lại gần mép sofa, cúi người nhặt chiếc kính râm vừa bị giẫm nứt lên rồi nheo mắt săm soi.

Tim tôi đập thình thịch như trống trận. Thôi xong rồi\dots\ Quả này thì đời tôi coi như tàn thật rồi! Chuẩn bị tinh thần đón nhận cơn thịnh nộ hóa quỷ của Miko thôi\dots!

\dots

\dots

Thế nhưng, trái ngược hoàn toàn với nỗi sợ hãi tột độ của tôi, Miko chỉ nghiêng đầu ngắm nghía chiếc kính một hồi, rồi buông một câu nhận xét hết sức bình thản: ``Hmm\dots\ Cái kính này trông kỳ lạ thật ne.''

Ủa\dots?

Tôi trố mắt kinh ngạc. Bà ấy\dots\ hoàn toàn không hề nổi giận với tôi ư?!

Thậm chí Miko còn cầm chiếc kính giơ lên cao trước mặt mọi người, hào hứng hỏi: ``Kuon-senpai, kính này của anh để quên à ne?''

Anh Kuon lắc đầu dứt khoát: ``Không phải. Mấy đứa thừa biết là anh có bao giờ đeo kính râm đâu. Sui-chan, cái này là của em mang tới à?''

Tôi nuốt nước bọt cái ực, xua tay lia lịa: ``K-Không phải của em đâu ạ!''

Game cất giọng thắc mắc: ``Lạ thật đấy. Thế rốt cuộc món đồ này từ đâu chui ra vậy? Hay là đồ đạc của chính cô đấy, Sakura-san?''

Miko bĩu môi giãy nảy: ``Làm gì có chuyện đó chứ ne! Tôi làm sao mà sở hữu một cặp kính ngầu lòi và cổ lỗ sĩ như thế này được!''

Thế là cả ba người họ xúm lại chụm đầu tranh luận xem rốt cuộc chủ nhân của chiếc kính râm kỳ lạ kia là ai, nhưng tuyệt nhiên chẳng có lấy một người nhận nó là của mình.

Quan trọng nhất là\dots\ toàn bộ những phản ứng và diễn biến đang diễn ra trước mắt tôi lúc này hoàn toàn trái ngược một trăm tám mươi độ so với những gì xảy ra trong cơn ác mộng kinh hoàng ban nãy.

Ít nhất thì\dots\ không có ai đùng đùng nổi giận, không có hắc khí tím ngùn ngụt, và cũng chẳng có ai nhảy bổ vào đòi ăn tươi nuốt sống tôi cả\dots\ nhỉ?

Kuon tặc lưỡi tiếc rẻ: ``Dù sao đi nữa thì món đồ này cũng hỏng mất rồi. Nứt toác cả mặt kính thế kia mà.''

Game bình phẩm: ``Thiết kế trông cũng khá bắt mắt đấy chứ\dots\ lại mang đậm phong cách hoài niệm nữa.''

Miko gật gù: ``Chắc chắn là loại này chỉ được bày bán vào đúng đợt năm 2002 thôi ne\dots\ Tiếc thật đấy\~{}''

Anh Tổng nheo mắt nhìn kỹ: ``\dots Mỗi tội nhìn hai con số hai bên chẳng đối xứng gì cả, trông cứ cấn cấn thế nào ấy.''

Game đồng tình: ``Công nhận thật.''

\dots Thậm chí đến nước này thì họ còn mải mê bàn luận về chiếc kính mà ngó lơ luôn cả sự tồn tại của tôi. Thú thực là lúc này tôi chẳng biết mình nên thở phào nhẹ nhõm vì thoát nạn, hay là nên cảm thấy hụt hẫng nữa.

Nhận thấy tôi cứ đứng chôn chân một chỗ im thin thít từ nãy tới giờ, anh Tổng bỗng quay sang nhìn tôi, ánh mắt lộ rõ vẻ ân cần quan tâm: ``Sui-chan, trông sắc mặt em kém lắm đấy. Em thấy trong người không được khỏe à?''

Tôi giật mình ngượng nghịu, vội xua tay lia lịa: ``D-Dạ em ổn mà! M-Mọi người không cần phải lo lắng cho em đâu ạ!'' Tôi cố nặn ra một nụ cười tươi tắn nhất có thể, dẫu cho mồ hôi lạnh vẫn đang rịn ra từng giọt trên thái dương.

Tất nhiên, với một người tinh tế và nhạy bén như Kuon-senpai, làm sao mà tôi có thể dễ dàng qua mặt anh ấy bằng nụ cười gượng gạo đó được. Vừa nghe tôi dứt lời, anh đã lập tức quay sang nhìn cô bạn thân của tôi, giọng nói hạ xuống đầy nghiêm túc: ``Miko-chan.''

Miko ngước mắt lên: ``Sao thế, ne?''

Kuon căn dặn: ``Em đưa con bé về nhà nghỉ ngơi giúp anh nhé. Để Sui-chan tự đi về một mình trong trạng thái thất thần thế này anh thấy không yên tâm chút nào đâu.''

Tôi vội xua tay từ chối: ``K-Không cần phiền phức thế đâu anh! Em thật sự tự đi về được mà\dots''

Nhưng Miko đã nhanh nhảu nắm chặt lấy cổ tay tôi, nở một nụ cười toe toét: ``Được rồi ne, để tôi hộ tống bà về tận cửa nhà luôn ne\~{}!''

Tôi méo mặt: ``Thôi nào Miko-chan, tôi lớn rồi tự đi được mà\dots\ Này, bà có thèm nghe tôi nói không đấy?!''

Kuon mỉm cười gật đầu: ``Cứ nghe lời anh đi. Hai đứa cứ về trước đi, việc ở công ty để anh báo lại với A-san một tiếng là được.''

Thế là, mặc cho tôi ra sức giãy giũa phân bua, Miko vẫn kiên quyết lôi tuột tôi ra khỏi cửa văn phòng. Từ phía sau lưng, vị trợ lý ảo Game vẫn không quên với theo một câu trêu chọc đầy thâm thúy: ``Hoshimachi-san với Đại PON-san đi đường cẩn thận nhé! Nhớ để mắt trông chừng nhau, đừng có ai lại bày trò gì đấy!''

Nàng Vu nữ lập tức quay đầu lại cãi nhem nhẻm: ``Tôi không có PON ne! Mồ! Ông nghĩ tôi là ai chứ hả?!'' Giọng cô ấy the thé đầy vẻ bực bội, nhưng ai cũng hiểu đó chỉ là sự phản kháng đáng yêu trước một câu đùa đã quá đỗi quen thuộc.

Rốt cuộc thì\dots\ danh tính chủ nhân thực sự của chiếc kính râm mang con số 2002 bí ẩn kia vẫn mãi mãi là một dấu chấm hỏi không lời đáp.

Thế nhưng, có một điều mà tôi hoàn toàn có thể thở phào chắc chắn: Ngày hôm nay, sẽ tuyệt đối không có bất kỳ ai bị trừng phạt, bị truy sát hay bị hãm hại cả.

\dots Ít nhất thì, tôi thành tâm hy vọng là như vậy!

\end{document}